\section{Diskussion \& Fazit}\label{sec:Diskussion}
In diesem Kaptiel soll die gesamte Arbeit reflektiert werden. In \cref{sec:Testbewertung} wurden bereits die Anforderungen mit
der entsprechenden Bilanz dargestellt. Dabei wurden nicht alle Anforderungen entsprechend optimal erfüllt. Aus diesem Grund
werden im folgenden Optimierungsschirtte dargelegt, die die auftretenden Fehler minimieren würden.
%
%
\subsection{Optimierungspotenzial}\label{sub:Optimierungspotezial}
Es wird in fünf verschiedenen Themen unterscheiden. Bei dem Aufbau einer solchen Analge bedarf es einen intensiven Kreilauf zwischen
Systemtest und Optimierung. Aus diesem Grund werden hier die von obern nach untern relevantesten Optimierungsschritte vorgeschlagen.
Wichtig zu benennen ist, dass die auftretenden Probleme zunächst keine schwäche der Arbeit sind, sondern bei richtiger 
Handhabung und dem beschrieben Kreislauf zu einer professionellen Herangehensweise werden.
%
\subsubsection{Geometrieabhängigkeit der Vereinzelung}\label{sub:HF-Vereinzelung}
Wie in der Messreihe T-05 (vgl. \cref{sec:VR_Vereinzelungstreue}) herausgefunden müssen die Sperrklinken/Endeffektoren auf die Hairpin-Geometrie
angpasst werden. Um eine Reduktion der Fehlerquellen zu ermöglcihen, gibt es dafür zwei mögliche Lösungsschritte. Die erste 
herangehensweise beinhaltet weitere Messreihen und eine neu Konstruktion für die Ermittlung der optimalen Endeffektorkombination.
Eine anderer für den allgemeinen Prozess sichereren Weg wäre die Erstellung einzelner Endeffektorkombinationen für die verwendete
Hairpin-Geometrie. Letzers könnte über ein generatives Design Modell entstehen, welches dann über ein AM-Drucker in kurzer zeit
hergestellt werden könnte. Den Wechsel könnte das Bedienpersonal dann durch das Lösen zweier Schrauben unkompliziert und schnell
druchführen. Eine nachfolgende Kalibierung des System ist ebenfalls einfach möglich. 
Diese Möglichkeiten würden Probleme stark dezimieren und der Mehraufwand ist nur bei einem Hairpin-Geometrie Wechsel nötig.

\subsubsection{Prozesssicherheit der Überführung}\label{sub:HF-Ueberfuehrung}


\subsubsection{Prozessueberwachung und Taktzeit}\label{sub:HF-Sensorik}


\subsubsection{Anbindung an die Zelle und Bedienung}\label{sub:HF-Anbindung}


\subsubsection{Werkstoffe und Verschleiss}\label{sub:HF-Werkstoffe}


\subsubsection{Alltagsbetrieb \& Limitationen}\label{sub:Erkenntnisse}
- Wo ist das System begrenzt?\\
- Was hätte anders gemacht werden müssen?\\


\subsection{Fazit}\label{sec:Fazit}
- Wurde die Forschungsfrage beantwortet?\\
- kritische Reflexion!\\
%
%
